% !TEX root = ../../../main.tex
% Block card: detection.dehshibi_spikes -- "Spike detection (Dehshibi, monolithic -- deprecated)".
\cardfield{What it does}
Runs the whole spike detector of \citet{dehshibi2021electrical} in one step: the same three
stages as the \texttt{dehshibi\_spikes} template, without showing you what happens in between.

\cardfield{When to use it}
Only to rerun an old recipe that names it. For new work use the template:
\blockref{preprocessing.wavelet_transform}, then \blockref{detection.wavelet_summation}, then
\blockref{detection.summation_threshold}. It gives the same spikes at the defaults and lets you
see each stage. That is why this block is filed under control, not detect.

\cardfield{What goes in, what comes out}
In goes a span of one channel; out come the spikes, one stretch each.

\cardfield{Settings that matter}
\emph{n\_p} is the shortest gap, in samples, between the rises and dips the method looks for
(60, which is 60\,s at 1\,Hz but only 6\,s at 10\,Hz). \emph{epsilon\_factor} is how prominent
they must be. The other two settings do nothing; they are kept so old recipes still load.

\cardfield{Watch out for}
Unlike the template, it has no limit on span length and no window setting (always 3000\,s), and
its spikes carry no labels or scores. It is as unselective on real recordings as the template
(see \blockref{detection.summation_threshold}).
